% cap3_metodologie.tex
\chapter{Metodologie și arhitectura sistemului}
\label{ch:metodologie}

Capitolul de față descrie principiile de proiectare care au ghidat construcția sistemului,
arhitectura de ansamblu pe cinci straturi, cerințele funcționale și nefuncționale, fluxul
de date bazat pe \gls{mqtt}, modelul relațional al bazei de date, alegerile tehnologice
justificate și, în final, metodologia contribuției originale — pipeline-ul de detectare
a anomaliilor comportamentale cu algoritmul Isolation Forest.

% ─────────────────────────────────────────────────────────────────────────────
\section{Principii de proiectare}
\label{sec:principii}
% ─────────────────────────────────────────────────────────────────────────────

Arhitectura sistemului a fost concepută respectând trei principii fundamentale care
influențează fiecare decizie tehnologică.

\textbf{Local-first.} Kleppmann et al.\ definesc paradigma \emph{local-first} ca
proiectarea în care toate datele utilizatorului sunt stocate și procesate pe dispozitivul
sau infrastructura locală, conexiunea la cloud reprezentând o opțiune auxiliară și nu o
cerință de funcționare~\cite{kleppmann2019localfirst}. Sistemul implementat respectă
integral această paradigmă: senzorarea, stocarea în \gls{iot}, inferența \gls{ml} și
controlul actuatoarelor se desfășoară exclusiv în rețeaua locală (\gls{lan}), fără niciun
trafic spre servicii externe. Consecința directă este că sistemul continuă să funcționeze
normal chiar dacă conexiunea la Internet este absentă.

\textbf{Privacy by design.} Principiile \emph{privacy by design}, consolidate în
Regulamentul European \gls{gdpr}~\cite{gdpr2016}, impun ca confidențialitatea să fie
integrată în arhitectură de la etapa de proiectare, și nu adăugată ulterior. Datele de
senzori — temperaturi, prezență, niveluri de gaz — sunt date cu caracter personal în
sensul GDPR atunci când permit identificarea activității ocupanților. Prin stocarea
exclusiv locală, sistemul elimină riscul de scurgere de date prin furnizori de cloud
terți, respectând principiul minimizării transferului de date.

\textbf{Edge intelligence.} Shi et al.\ propun modelul de calcul \emph{edge computing}
ca alternativă la centralizarea în cloud, prin deplasarea puterii de calcul cât mai aproape
de sursa datelor~\cite{shi2016edge}. În sistemul de față, motorul \gls{ml} (Isolation Forest)
rulează pe Raspberry Pi 3B+ — dispozitivul gateway — și nu pe un server cloud, obținând
latență de inferență de ordinul milisecundelor și independență operațională completă față
de Internet.

% ─────────────────────────────────────────────────────────────────────────────
\section{Arhitectura de ansamblu}
\label{sec:arhitectura}
% ─────────────────────────────────────────────────────────────────────────────

Sistemul este structurat în cinci straturi ierarhice, de la hardware la interfața
utilizator. \figref{fig:arhitectura} prezintă arhitectura de ansamblu; paragrafele
de mai jos descriu fiecare strat.

\begin{figure}[H]
  \centering
  \includegraphics[width=\singlefigure]{\pics/fig-arhitectura}
  \caption{Arhitectura de ansamblu a sistemului, organizată pe cinci straturi.
           Toate componentele operează exclusiv în rețeaua locală, fără dependență de cloud.}
  \label{fig:arhitectura}
\end{figure}

\textbf{Stratul 1 — Senzori și actuatori.}
Stratul fizic cuprinde senzorii de mediu și actuatoarele conectate la nodurile ESP32.
Senzorii măsoară: temperatura și umiditatea (2× DHT11), concentrația de gaze combustibile
(MQ-2, ieșire analogică), nivelul de luminozitate (2× \gls{ldr}) și prezența umană (un
singur senzor PIR, montat pe camera exterioară). Actuatoarele includ: 3 \gls{led}-uri
de iluminat, un ventilator de 5V, un servomotor SG90 (pentru controlul draperiilor) și
un buzzer activ pentru alerte sonore locale. Camera OV2640 (montată pe ESP32-CAM)
furnizează flux video MJPEG.

\textbf{Stratul 2 — Noduri ESP32.}
Două plăci de dezvoltare gestionează interfața cu hardware-ul fizic:
\begin{itemize}
  \item \emph{Nodul interior} (ESP32 WROOM-32) execută firmware Arduino C++ și comunică
    prin \gls{gpio}, \gls{adc} (10 biți, 0–1023) și \gls{pwm} cu toți senzorii și actuatoarele
    de interior.
  \item \emph{Nodul exterior} (ESP32-CAM, AI-Thinker, cu senzor de imagine OV2640) detectează
    mișcarea prin PIR pe GPIO13 și oferă server HTTP pe portul 80 pentru streaming video
    (\texttt{GET /stream}) și captură JPEG (\texttt{GET /capture}).
\end{itemize}
Ambele noduri comunică cu gateway-ul prin \gls{mqtt} peste Wi-Fi; camera este descoperită
de aplicația mobilă prin \gls{mdns} (\texttt{smarthome-cam.local}).

\textbf{Stratul 3 — Gateway Raspberry Pi.}
Raspberry Pi 3B+ îndeplinește rolul central al sistemului. Găzduiește:
\begin{itemize}
  \item brokerul Mosquitto (\gls{mqtt} 1883, autentificat), care primește și rutează
    toate mesajele publicate de noduri;
  \item baza de date PostgreSQL cu extensia TimescaleDB, pentru stocarea eficientă
    a seriilor temporale de citiri;
  \item motorul \gls{ml} Python (Isolation Forest), care antrenează și rulează inferența
    de detectare a anomaliilor direct pe dispozitivul edge.
\end{itemize}

\textbf{Stratul 4 — \gls{api} NestJS.}
Serviciul backend, scris în TypeScript cu framework-ul NestJS, expune un \gls{api} \gls{rest}
și un server WebSocket (Socket.IO) pe portul 3000. Comunicarea cu brokerul \gls{mqtt} se
realizează prin modulul \texttt{mqtt.service.ts}, care abonează NestJS la topicurile
de senzori și publică comenzi spre noduri. Autentificarea este realizată cu \gls{jwt}
și parole \emph{hash}-uite cu bcrypt.

\textbf{Stratul 5 — Aplicație mobilă.}
Interfața utilizatorului este o aplicație React Native (Expo SDK 54, TypeScript) ce
rulează pe Android și iOS. Primește date în timp real prin Socket.IO, afișează
valorile senzorilor, starea alertelor și permite controlul actuatoarelor. Descoperirea
automată a gateway-ului se face prin \gls{mdns} (\texttt{smarthome.local}).

% ─────────────────────────────────────────────────────────────────────────────
\section{Cerințe funcționale și nefuncționale}
\label{sec:cerinte}
% ─────────────────────────────────────────────────────────────────────────────

\subsection{Cerințe funcționale}

Cerințele funcționale au fost identificate pornind de la scenariile de utilizare reale
ale unui locatar. \figref{fig:usecase} prezintă diagrama use-case a sistemului.

\begin{figure}[H]
  \centering
  \includegraphics[width=\singlefigure]{\pics/fig-usecase}
  \caption{Diagrama use-case a sistemului smart home. Actorul principal este utilizatorul,
           iar motorul ML participă ca actor secundar (sistem) prin generarea automată
           de alerte de anomalie.}
  \label{fig:usecase}
\end{figure}

Principalele cerințe funcționale sunt:

\begin{enumerate}
  \item \textbf{Monitorizare senzori în timp real} — sistemul colectează și afișează
    temperaturi (2 camere), umiditate (2 camere), nivel de gaz, luminozitate (2 senzori)
    și prezență, cu actualizare la fiecare 5 secunde.
  \item \textbf{Control actuatori de la distanță} — utilizatorul poate acționa de pe
    telefon: LED-urile (pornit/oprit/toggle individual), ventilatorul, servomotorul
    draperiilor (0°/90°/180°), buzzer-ul și comanda \texttt{all\_off}.
  \item \textbf{Sistem de alerte} — alertele de tip \texttt{GAS\_DETECTED},
    \texttt{MOTION\_NIGHT}, \texttt{ALERT\_HEAT}, \texttt{ALERT\_COLD} și
    \texttt{ML\_ANOMALY} sunt stocate în baza de date și notificate în timp real
    prin WebSocket.
  \item \textbf{Vizualizare istoric} — aplicația permite consultarea istoricului
    grafic al citirilor pe intervale configurabile.
  \item \textbf{Flux video exterior} — aplicația afișează stream MJPEG live și
    captură snapshot de la camera exterioară.
  \item \textbf{Detectare anomalii comportamentale} — motorul \gls{ml} local
    generează alerte de tip \texttt{ML\_ANOMALY} pentru comportamente neobișnuite
    față de pattern-ul zilnic al locatarului.
  \item \textbf{Configurare Wi-Fi (provisioning)} — nodul poate primi noile
    credențiale Wi-Fi prin comanda \gls{mqtt} \texttt{wifi\_update} sau, ca
    fallback, prin hotspot-ul \texttt{SmartHome-Setup}.
\end{enumerate}

\subsection{Cerințe nefuncționale}

\begin{itemize}
  \item \textbf{Disponibilitate locală} — sistemul funcționează integral fără
    conexiune la Internet; nu există dependențe de servicii cloud externe.
  \item \textbf{Latență redusă} — de la publicarea datelor de senzori până la
    afișarea în aplicație, latența end-to-end vizată este sub 2 secunde.
  \item \textbf{Confidențialitate} — nicio dată de senzori sau video nu
    părăsește rețeaua locală.
  \item \textbf{Extensibilitate} — arhitectura permite adăugarea de noi noduri
    ESP32 fără modificări în backend, prin convenția de topicuri
    \texttt{smarthome\slb{} \{nodeId\}\slb{}sensors}.
  \item \textbf{Reziliență} — nodurile ESP32 implementează reconectare automată
    la Wi-Fi și \gls{mqtt}; LWT asigură notificarea automată a deconectărilor.
\end{itemize}

% ─────────────────────────────────────────────────────────────────────────────
\section{Fluxul de date și protocolul MQTT}
\label{sec:flux-date}
% ─────────────────────────────────────────────────────────────────────────────

\gls{mqtt} este protocolul de transport ales pentru comunicarea dintre noduri
și gateway. Standardul \gls{oasis} MQTT 3.1.1~\cite{oasis2019mqtt} definește un
protocol publish/subscribe binar, cu overhead minim, potrivit pentru dispozitive
cu resurse limitate~\cite{naik2017mqttcomparison}. Brokerul utilizat este
Mosquitto~\cite{light2017mosquitto}, rulând cu autentificare bazată pe parolă
pe portul 1883.

\figref{fig:flux-date} prezintă fluxul complet de date din sistem.

\begin{figure}[H]
  \centering
  \includegraphics[width=\singlefigure]{\pics/fig-flux-date}
  \caption{Fluxul de date în sistem. Datele de senzori circulă de la ESP32 prin
           Mosquitto spre NestJS, care le persistă în TimescaleDB, le transmite
           aplicației prin WebSocket și le evaluează prin motorul ML.
           Comenzile urmează traseul invers (linie punctată).}
  \label{fig:flux-date}
\end{figure}

\subsection{Topicuri MQTT și contractul de date}

\tabref{tab:mqtt-topics} detaliază topicurile \gls{mqtt} utilizate, direcția
fluxului, nivelul QoS și rolul fiecăruia.

\begin{table}[H]
  \centering
  \caption{Topicurile MQTT ale sistemului smart home.}
  \label{tab:mqtt-topics}
  \small
  \begin{tabular}{@{}p{4.4cm} p{2.4cm} p{1.5cm} p{4.6cm}@{}}
    \toprule
    \textbf{Topic} & \textbf{Direcție} & \textbf{QoS} & \textbf{Descriere} \\
    \midrule
    \texttt{smarthome\slb{} \{nodeId\}\slb{}sensors}
      & ESP32 → broker
      & 0
      & Citiri periodice (la 5 s) în format JSON; payload conține
        temperaturi, umiditate, gaz, luminozitate. \\
    \addlinespace
    \texttt{smarthome\slb{} \{nodeId\}\slb{}status}
      & ESP32 → broker
      & 1, retained
      & Starea nodului: \texttt{online} la conectare,
        \texttt{offline} prin Last Will (LWT). \\
    \addlinespace
    \texttt{smarthome\slb{} \{nodeId\}\slb{}commands}
      & broker → ESP32
      & 1
      & Comenzi pentru actuatori: \texttt{led\_on}, \texttt{led\_off},
        \texttt{led\_toggle}, \texttt{fan\_on}, \texttt{fan\_off},
        \texttt{servo\_move}, \texttt{buzzer\_beep}, \texttt{all\_off},
        \texttt{wifi\_update}. \\
    \addlinespace
    \texttt{smarthome/alerts}
      & ESP32 / ML → broker
      & 1
      & Alerte de urgență: gaz detectat (de pe nod) sau anomalie ML
        (de pe gateway). \\
    \bottomrule
  \end{tabular}
\end{table}

Payloadul publicat de nodul interior pe topicul de senzori este un obiect JSON
cu câmpurile prezentate în fragmentul de mai jos (contract comun între firmware,
backend și aplicația mobilă):

\begin{lstlisting}[language=TypeScript, caption={Payload MQTT publicat de nodul interior la fiecare 5 secunde.}]
{
  "nodeId":       "esp32_node_a",
  "location":     "interior",
  "temperature":  22.5,   // DHT11 #1, living (°C)
  "humidity":     55,     // DHT11 #1, living (%RH)
  "temperature2": 21.8,   // DHT11 #2, dormitor (°C)
  "humidity2":    58,     // DHT11 #2, dormitor (%RH)
  "gasLevel":     245,    // MQ-2 ADC 0-1023
  "gasAlert":     false,  // true dacă gasLevel > 350
  "light1":       410,    // LDR #1 ADC 0-1023
  "light2":       120,    // LDR #2 ADC 0-1023
  "lightLux":     410     // alias light1 (compatibilitate)
}
\end{lstlisting}

\subsection{Last Will and Testament și QoS}

La conectare, nodul configurează mesajul \emph{Last Will and Testament} (LWT) pe topicul
\texttt{smarthome\slb{} \{nodeId\}\slb{}status} cu payload \texttt{\{"status":"offline"\}} și flag
\emph{retained}. Dacă conexiunea TCP se întrerupe fără deconectare explicită, brokerul
publică automat mesajul LWT, permițând NestJS să marcheze nodul ca offline în baza
de date și să notifice aplicația mobilă prin WebSocket.

Datele de senzori se publică cu QoS~0 (fire-and-forget), acceptabil pentru telemetrie
continuă cu cadență mare (pierderea unui pachet din 720 de zi este neglijabilă).
Comenzile și alertele folosesc QoS~1 (cel puțin o dată), asigurând livrarea
confirmată.

\subsection{Descoperire prin mDNS și provisioning Wi-Fi}

Descoperirea automată a broker-ului și a camerei utilizează \gls{mdns} (RFC 6762~\cite{rfc6762mdns}):
ESP32 rezolvă \texttt{smarthome.local} la adresa IP curentă a Raspberry Pi, fără
configurare manuală. Camera exterioară este disponibilă la \texttt{smarthome-cam.local}.

Provisioningul Wi-Fi urmează o schemă cu fallback: credențialele curente (SSID, parolă,
adresa broker) sunt persistate în memoria NVS (\emph{Non-Volatile Storage}) a ESP32.
La pornire, nodul încearcă conectarea la rețeaua salvată; în caz de eșec, revine automat
la hotspot-ul de configurare \texttt{SmartHome-Setup} (IP fix \texttt{192.168.4.1}).
Credențialele pot fi actualizate în orice moment prin comanda \gls{mqtt}
\texttt{wifi\_update \{ssid, password, mqtt\_host\}}, după care nodul repornește
și se reconectează la noua rețea.

% ─────────────────────────────────────────────────────────────────────────────
\section{Modelul de date}
\label{sec:model-date}
% ─────────────────────────────────────────────────────────────────────────────

Baza de date rulează pe PostgreSQL~\cite{postgresql2024} cu extensia
TimescaleDB~\cite{timescaledb2017}, specializată în stocarea și interogarea eficientă
a seriilor temporale~\cite{jensen2017timeseries}. \figref{fig:er} prezintă diagrama
entitate-relație.

\begin{figure}[H]
  \centering
  \includegraphics[width=\singlefigure]{\pics/fig-er}
  \caption{Diagrama entitate-relație a bazei de date. Tabelul \texttt{sensor\_readings}
           este hypertable TimescaleDB particionat automat pe coloana \texttt{time}.}
  \label{fig:er}
\end{figure}

\textbf{Tabelul \texttt{sensor\_readings} (hypertable).}
Este tabelul central al sistemului, convertit în hypertable TimescaleDB cu partiționare
automată pe coloana \texttt{time}. Stochează fiecare citire publicată de un nod ESP32:

\begin{itemize}
  \item \texttt{time} — \texttt{TIMESTAMPTZ}, coloana de partiționare;
  \item \texttt{node\_id}, \texttt{location} — identificatorul și locația nodului;
  \item \texttt{temperature}, \texttt{humidity} — citirile DHT11~\#1 (living);
  \item \texttt{temperature2}, \texttt{humidity2} — citirile DHT11~\#2 (dormitor),
    \texttt{NULLABLE} (câmpuri absente în citiri mai vechi sau în caz de eroare senzor);
  \item \texttt{gas\_level}, \texttt{gas\_alert} — valoarea \gls{adc} MQ-2 și indicatorul
    de alertă;
  \item \texttt{light1}, \texttt{light2}, \texttt{light\_lux} — valorile celor două
    \gls{ldr}-uri și alias-ul de compatibilitate.
\end{itemize}

TimescaleDB comprimă automat chunk-urile mai vechi de 7 zile și permite interogări
de agregare temporală cu funcții native (\texttt{time\_bucket}, \texttt{first},
\texttt{last}) considerabil mai rapide decât PostgreSQL simplu pe seturi mari
de date~\cite{timescaledb2017}.

\textbf{Tabelul \texttt{alerts}.}
Stochează alertele generate, cu câmpurile: \texttt{id}, \texttt{time}, \texttt{node\_id},
\texttt{alert\_type} (unul dintre \texttt{GAS\_DETECTED}, \texttt{MOTION\_NIGHT},
\texttt{ALERT\_HEAT}, \texttt{ALERT\_COLD}, \texttt{ML\_ANOMALY}),
\texttt{severity}, \texttt{location}, \texttt{details} (JSON cu context),
\texttt{acknowledged} (boolean).

\textbf{Tabelul \texttt{nodes}.}
Registrul nodurilor active: \texttt{id}, \texttt{node\_id}, \texttt{location},
\texttt{last\_seen}, \texttt{online}, \texttt{node\_type}.

\textbf{Tabelul \texttt{automation\_rules}.}
Permite definirea de reguli de automatizare configurabile: senzor, operator,
prag (\texttt{threshold}), interval orar (\texttt{time\_from}, \texttt{time\_to}),
tip de acțiune (\texttt{action\_type}), stare activă.

\textbf{Tabelele \texttt{users} și \texttt{ml\_models}.}
\texttt{users} păstrează credențialele autentificate (username, \texttt{password\_hash}
produs prin bcrypt). \texttt{ml\_models} înregistrează metadatele fiecărui model antrenat:
versiune (formatul \texttt{vAAAAMMZZ\_HHMM}), numărul de rânduri de antrenare,
valoarea \texttt{contamination}, lista feature-urilor și câmpul \texttt{is\_active}
care identifică modelul curent utilizat în inferență.

% ─────────────────────────────────────────────────────────────────────────────
\section{Alegeri tehnologice justificate}
\label{sec:tehnologii}
% ─────────────────────────────────────────────────────────────────────────────

\tabref{tab:tehnologii} sintetizează principalele alegeri tehnologice și justificarea
lor în raport cu cerințele sistemului.

\begin{table}[H]
  \centering
  \caption{Alegeri tehnologice — cerință, soluție adoptată și justificare.}
  \label{tab:tehnologii}
  \begin{tabular}{@{}p{3.5cm} p{3.2cm} p{6.5cm}@{}}
    \toprule
    \textbf{Cerință} & \textbf{Tehnologie} & \textbf{Justificare} \\
    \midrule
    Microcontroler \gls{iot} cu Wi-Fi
      & ESP32 WROOM-32~\cite{espressif2023esp32}
      & Dual-core 240~MHz, 520~KB SRAM, Wi-Fi 802.11b/g/n integrat, suport Arduino;
        cost redus (\textasciitilde5~USD); \gls{adc} 12-biți pe mai mulți pini. \\
    \addlinespace
    Transport mesaje pub/sub
      & MQTT 3.1.1 / Mosquitto~\cite{oasis2019mqtt,light2017mosquitto}
      & Protocol binar cu overhead minim (\textasciitilde2~B header); QoS 0–2;
        LWT pentru detectarea deconectărilor; potrivit pentru rețele locale cu
        lățime de bandă redusă. \\
    \addlinespace
    Stocare serii temporale
      & PostgreSQL + TimescaleDB~\cite{timescaledb2017,postgresql2024}
      & Partiționare automată (hypertable) pe timp; compresie nativă; funcții de
        agregare temporală optimizate; interoperabilitate SQL deplină cu TypeORM. \\
    \addlinespace
    Backend \gls{api}
      & NestJS + TypeScript~\cite{nestjs2024}
      & Arhitectură modulară cu injecție de dependențe; integrare nativă WebSocket
        (Socket.IO); TypeORM pentru acces la baza de date; suport \gls{jwt} din cutie. \\
    \addlinespace
    Aplicație mobilă
      & React Native / Expo~\cite{reactnative2024,expo2024}
      & Cod TypeScript partajat între iOS și Android; Expo SDK 54 simplifică
        build-ul; ecosistem bogat (React Query, Zustand, Socket.IO client). \\
    \addlinespace
    Detectare anomalii
      & Isolation Forest~\cite{liu2008isolation,liu2012isolationtkdd}
      & Algoritm nesupervizat, fără etichete; complexitate de antrenare
        $O(n \log n)$; consum de memorie redus (\textasciitilde4~MB model pkl);
        implementare matură în scikit-learn~\cite{pedregosa2011scikit};
        rulabil pe Raspberry Pi 3B+ fără accelerator hardware. \\
    \addlinespace
    Autentificare
      & JWT + bcrypt~\cite{jones2015jwt,provos2003bcrypt}
      & Tokenuri fără stare, ușor de validat la fiecare request \gls{rest};
        bcrypt cu factor de cost adaptabil pentru hash-uirea parolelor. \\
    \bottomrule
  \end{tabular}
\end{table}

% ─────────────────────────────────────────────────────────────────────────────
\section{Pipeline-ul de detectare a anomaliilor — contribuția originală}
\label{sec:ml-pipeline}
% ─────────────────────────────────────────────────────────────────────────────

Contribuția originală centrală a lucrării este un pipeline de detectare a anomaliilor
comportamentale bazat pe algoritmul Isolation Forest, implementat în Python și rulând
local pe Raspberry Pi 3B+. Spre deosebire de sistemele care transmit datele la cloud
pentru analiză, abordarea propusă menține toate datele și calculele în rețeaua locală,
respectând principiile de confidențialitate enunțate în \secref{sec:principii}.

\figref{fig:ml-pipeline} prezintă structura completă a pipeline-ului.

\begin{figure}[H]
  \centering
  \includegraphics[width=\singlefigure]{\pics/fig-ml-pipeline}
  \caption{Pipeline-ul complet al motorului ML: de la colectarea datelor din
           TimescaleDB, prin ingineria trăsăturilor și normalizare, până la
           antrenarea periodică a modelului Isolation Forest și inferența în timp real.
           Bucla de reantrenare zilnică și mecanismul de cold-start sunt evidențiate.}
  \label{fig:ml-pipeline}
\end{figure}

\subsection{Algoritmul Isolation Forest}

Isolation Forest a fost propus de Liu, Ting și Zhou~\cite{liu2008isolation} și
extins ulterior~\cite{liu2012isolationtkdd} ca algoritm de detectare a anomaliilor
bazat pe proprietatea că punctele anomale sunt mai ușor de „izolat" prin partiționări
aleatoare ale spațiului de trăsături față de punctele normale. Formal, un arbore
de izolare (\emph{iTree}) partiționează recursiv un subsamplu de date prin selecție
aleatoare a unei trăsături și a unui prag de separare. Lungimea medie a căii de la
rădăcină până la un punct $\mathbf{x}$ în toți arborii constituie scorul de anomalie.

Scorul de anomalie $s(\mathbf{x}, n)$ este definit ca:

\begin{equation}
  s(\mathbf{x}, n) = 2^{-\dfrac{E[h(\mathbf{x})]}{c(n)}},
  \label{eq:if-score}
\end{equation}

unde $h(\mathbf{x})$ este lungimea căii parcurse de $\mathbf{x}$ în arbore, $E[\cdot]$
este media peste toți arborii din ansamblu, $n$ este dimensiunea subsamplului de antrenare
și $c(n) = 2H(n-1) - \tfrac{2(n-1)}{n}$ este lungimea medie a căii pentru $n$ noduri
($H$ — număr armonic). Un scor $s \to 1$ indică o anomalie pronunțată (cale scurtă),
iar $s \to 0$ indică un punct normal. Implementarea din scikit-learn~\cite{pedregosa2011scikit}
returnează \texttt{score\_samples(x)} ca valori negative, cu semnatica: cu cât scorul
este mai negativ, cu atât punctul este mai anormal.

\subsection{Ingineria trăsăturilor (\texttt{features.py})}

Modulul \texttt{features.py} construiește 12 trăsături numerice din citirile brute ale
senzorilor, în ordinea fixă din \tabref{tab:features}. Ordinea este identică în
antrenare și în inferență — orice modificare necesită reantrenare completă a modelului.

\begin{table}[H]
  \centering
  \caption{Cele 12 trăsături ale modelului Isolation Forest, în ordinea fixă din \texttt{FEATURE\_NAMES}.}
  \label{tab:features}
  \begin{tabular}{@{}r l p{7.5cm}@{}}
    \toprule
    \textbf{\#} & \textbf{Trăsătură} & \textbf{Descriere și sursă} \\
    \midrule
    0  & \texttt{temperature}      & Temperatura DHT11~\#1 (living), °C. \\
    1  & \texttt{humidity}         & Umiditatea relativă DHT11~\#1 (\%RH). \\
    2  & \texttt{gas\_normalized}  & Valoarea ADC MQ-2 împărțită la 4095
                                     (normalizare în $[0,1]$). \\
    3  & \texttt{lux\_normalized}  & $\log_{1p}(\texttt{light1}) / 7{,}0$ —
                                     luminozitate LDR~\#1 normalizată logaritmic
                                     pe scala ADC 0–1023
                                     ($\log_{1p}(1023) \approx 6{,}93$). \\
    4  & \texttt{motion}           & Prezența umană detectorului PIR (0/1). \\
    5  & \texttt{hour\_of\_day}    & Ora din zi (0–23), captează ritmul circadian. \\
    6  & \texttt{day\_of\_week}    & Ziua săptămânii (0 = luni, 6 = duminică). \\
    7  & \texttt{is\_weekend}      & 1 dacă ziua este sâmbătă sau duminică. \\
    8  & \texttt{is\_night}        & 1 dacă ora $\geq 23$ sau ora $< 6$. \\
    9  & \texttt{temp\_delta\_1h}  & Diferența de temperatură față de citirile
                                     cu 2 eșantioane în urmă (aproximativ 10~s). \\
    10 & \texttt{motion\_count\_10}& Suma evenimentelor de mișcare în ultimele
                                     10 citiri (fereastră glisantă). \\
    11 & \texttt{temp\_avg}        & Media celor doi senzori DHT11; dacă DHT11~\#2
                                     lipsește, se folosește valoarea DHT11~\#1. \\
    \bottomrule
  \end{tabular}
\end{table}

\subsection{Normalizare}

Înainte de antrenare, trăsăturile sunt normalizate cu \texttt{StandardScaler} din
scikit-learn, care transformă fiecare trăsătură la medie 0 și deviație standard 1:
\begin{equation}
  z_i = \frac{x_i - \mu_i}{\sigma_i}.
  \label{eq:standardscaler}
\end{equation}
Parametrii $\mu_i$ și $\sigma_i$ sunt calculați pe datele de antrenare și salvați în
\texttt{scaler.pkl} alături de model. La inferență, aceeași transformare este aplicată
fiecărei noi citiri înainte de evaluarea scorului.

\subsection{Antrenarea modelului (\texttt{train.py})}

Modulul \texttt{train.py} implementează ciclul de antrenare:

\begin{enumerate}
  \item Citește din tabela \texttt{sensor\_readings} toate înregistrările din
    ultimele 30 de zile cu \texttt{gas\_alert = FALSE} (citeșterile din perioadele
    de alertă sunt excluse pentru a nu contamina distribuția de normalitate).
  \item Verifică pragul minim de 500 de rânduri; dacă datele sunt insuficiente,
    antrenarea este întreruptă cu mesaj informativ.
  \item Apelează \texttt{build\_features(df)} și normalizează cu \texttt{StandardScaler}.
  \item Antrenează un \texttt{IsolationForest} cu parametrii:
    \texttt{n\_estimators=100}, \texttt{contamination=0.05}, \texttt{random\_state=42},
    \texttt{n\_jobs=-1} (toate nucleele disponibile).
  \item Salvează \texttt{isolation\_forest.pkl} și \texttt{scaler.pkl} în directorul
    \texttt{models/}.
  \item Înregistrează metadatele modelului nou în tabela \texttt{ml\_models} și
    marchează toate versiunile anterioare ca inactive.
\end{enumerate}

Parametrul \texttt{contamination=0.05} afirmă că modelul estimează că 5\% din datele
de antrenare reprezintă comportamente atipice. Această valoare este conservatoare —
subestimarea duce la mai puține fals-negative în faza de deploy. Reantrenarea este
planificată zilnic la ora 03:00 prin \texttt{scheduler.py}, rulat ca serviciu systemd
\texttt{smarthome-ml.service} pe Raspberry Pi.

\subsection{Inferența în timp real (\texttt{infer.py})}

Clasa \texttt{AnomalyDetector} din \texttt{infer.py} este instanțiată la pornirea
serviciului NestJS și menținută în memorie. La fiecare citire nouă de senzori:

\begin{enumerate}
  \item Se construiește un DataFrame cu o singură linie din payload-ul JSON primit.
  \item Se aplică \texttt{build\_features} și \texttt{scaler.transform}.
  \item Se apelează \texttt{model.score\_samples(X\_scaled)}, care returnează un scalar
    real; notăm scorul ca $s \in \mathbb{R}$.
  \item Dacă $s < \tau$ (unde $\tau = -0{,}1$ este pragul \texttt{ANOMALY\_THRESHOLD}),
    citirea este clasificată drept anomalie și se creează o alertă de tip
    \texttt{ML\_ANOMALY} în baza de date, notificată imediat aplicației prin WebSocket.
\end{enumerate}

Metoda returnează un obiect cu câmpurile \texttt{is\_anomaly} (boolean), \texttt{score}
(valoare numerică), \texttt{threshold} și \texttt{model\_active}. Dacă fișierele
\texttt{.pkl} nu există, câmpul \texttt{model\_active} este \texttt{false} și
nicio alertă ML nu este generată.

\subsection{Mecanismul de cold-start}

În primele aproximativ 14 zile de funcționare a sistemului, nu există suficiente date
pentru a antrena modelul (minimul de 500 de rânduri echivalând cu circa 42 de minute
de date, dar o fereastră de 30 de zile este necesară pentru a surprinde variabilitatea
comportamentală zilnică și săptămânală). În această perioadă de cold-start, motorul
ML returnează \texttt{model\_active: false} și sistemul se bazează pe reguli de prag
hardcodate pentru generarea alertelor:

\begin{itemize}
  \item \texttt{gas\_level > 300} → alertă \texttt{GAS\_DETECTED};
  \item \texttt{temperature > 35°C} → alertă \texttt{ALERT\_HEAT};
  \item \texttt{temperature < 5°C} → alertă \texttt{ALERT\_COLD};
  \item mișcare detectată în intervalul 23:00–06:00 → alertă \texttt{MOTION\_NIGHT}.
\end{itemize}

Această strategie asigură că sistemul oferă protecție de bază chiar înainte ca modelul
\gls{ml} să fie disponibil.

% ─────────────────────────────────────────────────────────────────────────────
\section{Considerații de securitate}
\label{sec:securitate}
% ─────────────────────────────────────────────────────────────────────────────

Modelul de securitate al sistemului este construit pe trei niveluri complementare.

\textbf{Autentificare și autorizare.}
Accesul la \gls{api}-ul \gls{rest} și la conexiunile WebSocket este protejat prin
\gls{jwt}~\cite{jones2015jwt}. La autentificarea reușită, serverul emite un token semnat
cu o cheie secretă; toate request-urile ulterioare trebuie să includă tokenul în header-ul
\texttt{Authorization: Bearer <token>}. Parolele utilizatorilor sunt stocate exclusiv
ca hash-uri bcrypt~\cite{provos2003bcrypt} cu factor de cost adaptabil (implicit 10
runde), eliminând riscul expunerii parolelor în clar chiar în cazul unui acces
neautorizat la baza de date.

\textbf{Autentificare broker MQTT.}
Brokerul Mosquitto este configurat cu autentificare prin user/parolă (utilizatorul
\texttt{smarthome}, credențiale stocate în fișierul Mosquitto \texttt{passwd}),
permițând conectarea doar nodurilor și serviciilor cunoscute. Comunicarea se
desfășoară în rețeaua locală și nu este expusă direct la Internet.

\textbf{Izolare în rețeaua locală.}
Nicio componentă a sistemului nu expune porturi pe Internet. Gateway-ul Raspberry Pi,
nodurile ESP32 și aplicația mobilă comunică exclusiv în \gls{lan}. Această izolare
reprezintă o măsură de securitate perimetrală complementară celor aplicative,
conformă cu principiul \emph{privacy by design} din \gls{gdpr}~\cite{gdpr2016}.
Datele cu caracter personal (activitate comportamentală, prezență umană) nu
părăsesc niciodată rețeaua locală, eliminând riscurile asociate transferului
transfrontalier de date sau breșelor de securitate la nivel de furnizor cloud.

\textbf{Minimizarea suprafeței de atac.}
Firmware-ul ESP32 nu expune nicio interfață HTTP sau Telnet în mod normal de funcționare;
singurul vector de comunicare externă este conexiunea \gls{mqtt} autentificată.
ESP32-CAM expune server HTTP pe portul 80 exclusiv în rețeaua locală, accesul nefiind
publicat prin NAT sau port-forwarding.
